\documentclass[11pt]{article}
\usepackage[a4paper,margin=24mm]{geometry}
\usepackage{amsmath,amssymb,amsthm,booktabs,array,longtable,hyperref}
\hypersetup{colorlinks=true,linkcolor=blue,citecolor=blue,urlcolor=blue,
  pdftitle={Polynomial-flow extensions of Su's theorem on Bala's period conjectures},
  pdfauthor={Haobo Ma; Wenlin Zhang}}
\setlength{\emergencystretch}{2em}
\newtheorem{theorem}{Theorem}
\newtheorem{lemma}[theorem]{Lemma}
\newtheorem{corollary}[theorem]{Corollary}
\newtheorem{proposition}[theorem]{Proposition}
\newcommand{\Z}{\mathbb Z}
\newcommand{\Q}{\mathbb Q}
\newcommand{\F}{\mathbb F}
\newcommand{\ph}{\varphi}
\newcommand{\ev}{\operatorname{ev}_0}
\title{Polynomial-flow extensions of Su's theorem\\on Bala's period conjectures}
\author{Haobo Ma\\\small ChronoAI Pte Ltd, Singapore
  \and Wenlin Zhang\\\small National University of Singapore\\\small ChronoAI Pte Ltd, Singapore}
\date{14 November 2026}
\begin{document}
\maketitle
\begin{center}\small Licensed under Creative Commons Attribution 4.0 International
  (CC BY 4.0).\end{center}
\begin{abstract}
We give an alternative proof and an extension of Lien-Hung Su's theorem
(31 October 2026, Theorem 3.2 and Corollary 3.3, \cite{su}) resolving
Peter Bala's conjecture in OEIS A281180, Comments, revision 86,
signed 18 February 2026, and its four companion conjectures.
The proof begins with a uniform
operator identity for every integral polynomial vector field
$P(z)\,d/dz$ with $P(0)=1$.
For the quadratic flows $S'=C^a$, $C'=SC^{a-1}$, we extend the eventual
period $\ph(m)^2/2$ to every integer $a\ge1$, every integral polynomial
observable at even derivative indices, and every $m\ge3$.
We give explicit starting indices, prove residue-row and derivative
identities from the neighboring entries, and supply exact finite-matrix
certificates for the source examples.
\end{abstract}

\section{The question and the earlier proof}
Let $S\in\Q[[x]]$ be the unique formal solution
\begin{equation}\label{eq:source}
 S'=(1+S^2)^2,\qquad S(0)=0,\qquad
 a(n)=(2n-1)![x^{2n-1}]S(x)\quad(n\ge1).
\end{equation}
This is OEIS A281180 \cite{oeis180}. Bala's comment asks whether,
for every positive integer $m$, the sequence $a(n)$ is eventually
periodic modulo $m$, with period dividing $\ph(m)^2$.
Thus the assertion is a congruence with that prescribed shift, not
an assertion that this shift is the least period.

Su's Theorem 3.2 and Corollary 3.3 \cite{su} already prove this and
the corresponding assertions for A281181, A281183, A281184 and
A281440. They also give the stronger half-square shift for $m\ge3$,
with smaller starting-index bounds than the uniform bounds below.
Our purpose is an alternative proof through truncated polynomial
rings and a direct explanation of the distinction between the
derivative index and the sequence index, followed by an extension
from even exponents to all integral exponents. The correspondence between
inverse functions and iterated differential operators is classical;
Bala \cite{bala} gives the formal identity in the notation relevant
here, and Dominici \cite{dominici} treats nested-derivative expansions.
Danesh's weighted Stirling-transform criterion \cite[Theorems 3.1
and 4.1]{danesh} gives other eventual-period theorems through valuation
bounds and exponential sums. Here the finite objects are differential
operators on truncated polynomial rings.

For $m>1$ put
\begin{equation}\label{eq:bounds}
 M(m)=\max_{p^r\parallel m}p^{2r-1},\qquad
 N(m)=\left\lceil\frac{M(m)+1}{2}\right\rceil.
\end{equation}
\begin{theorem}\label{thm:main}
The integers $a(n)$ in \eqref{eq:source} satisfy, for every $m\ge3$,
\[
 a\bigl(n+\ph(m)^2/2\bigr)\equiv a(n)\pmod m
 \qquad(n\ge N(m)).
\]
Modulo $2$ they are zero for $n\ge2$; modulo $1$ the assertion is
immediate. Consequently the shift $\ph(m)^2$ is valid eventually
for every positive integer $m$.
\end{theorem}

\section{An operator identity for polynomial flows}
For $P,f\in\Z[z]$ with $P(0)=1$, write
$D=P(z)\,d/dz$ and $u_j=(D^jf)(0)$. These are integers.
The central statement is the following one, with no restriction
on the degree of $P$ or $f$.

\begin{lemma}\label{lem:operator}
For every prime $p$, $r\ge1$ and $j\ge p^{2r-1}$,
\begin{equation}\label{eq:operator}
 u_{j+\ph(p^r)^2}\equiv u_j\pmod{p^r}.
\end{equation}
The same conclusion modulo $m>1$ holds for $j\ge M(m)$
with shift $\ph(m)^2$.
\end{lemma}

\begin{proof}
First work in $A_1=\F_p[z]/(z^p)$. The derivative descends to
this ring, and $P$ is a unit since $P(0)=1$. The iterated Leibniz
rule shows that $D^p$ is again a derivation: the intermediate
binomial coefficients vanish in characteristic $p$. Every
derivation on $A_1$ is determined by its value on $z$, so
$D^p=\lambda D$ for some $\lambda\in A_1$.
The commutation of $D$ and $D^p$ implies
$D(\lambda)D(z)=0$. As $D(z)=P$ is a unit, $D(\lambda)=0$,
and then $\lambda'=0$. A representative of degree less than $p$
with derivative zero is constant. Thus $\lambda=\alpha\in\F_p$.

Set $t=(p-1)^2$. If $\alpha=0$, then $D^p=0$.
If $\alpha\ne0$, induction gives
$D^{1+v(p-1)}=\alpha^vD$ for $v\ge0$, so Fermat's theorem gives
$D^{1+t}=D$. Both cases imply
\begin{equation}\label{eq:firstjet}
 Q:=D^p(D^t-I)=0\quad\hbox{on }A_1.
\end{equation}
This includes $p=2$.

Now regard $D$ and $Q$ as operators on $\F_p[z]$.
Equation \eqref{eq:firstjet} means that $Q(g)$ is divisible by
$z^p$ for every polynomial $g$. Since $D(z^{ph})=0$ for every
$h\ge0$, the operator $Q$ commutes with multiplication by $z^{ph}$.
It follows inductively that $Q^h(g)$ is divisible by $z^{ph}$.
In particular, on $A_r=\F_p[z]/(z^{p^r})$, one has
$Q^{p^{r-1}}=0$. Frobenius in the commutative algebra generated
by $D$ therefore gives
\begin{equation}\label{eq:modp}
 D^{p^r}\bigl(D^{(p-1)^2p^{r-1}}-I\bigr)=0
 \quad\hbox{on }A_r.
\end{equation}

Let $L=p^r$, and let $B$ be the integer matrix of the operation
``apply $P\,d/dz$ and discard terms of degree at least $L$''
in the basis $1,z,\ldots,z^{L-1}$. Its reduction modulo $p$
satisfies \eqref{eq:modp}. For commuting integer matrices $U,V$,
the implication
\[
 U\equiv V\pmod p\quad\Longrightarrow\quad
 U^{p^{r-1}}\equiv V^{p^{r-1}}\pmod{p^r}
\]
follows by repeated binomial expansion. Indeed, if $U-V$ is
divisible by $p^s$, $s\ge1$, each term of $U^p-V^p$ is
divisible by $p^{s+1}$: the intermediate binomial coefficients
are divisible by $p$, and $sp\ge s+1$ for the last term.
Apply this implication to
$U=B^{L+(p-1)^2p^{r-1}}$ and $V=B^L$. We obtain
\begin{equation}\label{eq:lift}
 B^{p^{2r-1}}\bigl(B^{(p-1)^2p^{2r-2}}-I\bigr)
 \equiv0\pmod{p^r}.
\end{equation}
Modulo $p^r$, the ideal $(z^L)$ is stable under $P\,d/dz$,
since the derivative of $z^L$ is zero. Hence iterates of $B$
represent iterates of $D$ in the quotient, and evaluation at zero
gives \eqref{eq:operator}. Multiplying \eqref{eq:lift} by
further powers of $B$ proves it for all the indicated $j$.

Finally $\ph(p^r)\mid\ph(m)$ for every $p^r\parallel m$,
by the product formula for the totient. Thus each local shift
divides $\ph(m)^2$. Repeat a local congruence a whole number
of times, always at indices at least $M(m)$, and apply the
Chinese remainder theorem.
\end{proof}

\begin{proof}[Proof of Theorem~\ref{thm:main}]
For $P=(1+z^2)^2$ set
$\mathcal F(f)=\sum_{j\ge0}(D^jf)(0)x^j/j!$.
The iterated Leibniz rule proves
$\mathcal F(fg)=\mathcal F(f)\mathcal F(g)$, and
$\mathcal F(f)'=\mathcal F(Df)$.
Consequently $\mathcal F(z)$ satisfies \eqref{eq:source};
uniqueness follows by solving each next coefficient in its
formal differential equation. This proves both
$S=\mathcal F(z)$ and the integrality of $a(n)=u_{2n-1}$.
If $m\ge3$, then $\ph(m)$ is even. The shift $\ph(m)^2/2$
in $n$ is exactly the shift $\ph(m)^2$ in $2n-1$.
Since $2n-1\ge M(m)$ when $n\ge N(m)$,
Lemma~\ref{lem:operator} proves the claim. For $m=2$ the lemma
gives $u_{j+1}=u_j$ for $j\ge2$; directly $u_2=0$,
so $a(n)=0$ for $n\ge2$.
\end{proof}

\section{All integral exponents and the companion conjectures}
Let $a\ge1$ and use the integral quadratic algebra
\[
 R=\Z[s,c]/(c^2-s^2-1),\qquad
 D_a(s)=c^a,\quad D_a(c)=sc^{a-1},\quad E(f)=f(0,1).
\]
The derivation preserves the defining relation. Its exponential
flow gives series $S_a,C_a$ with
\[
 S_a'=C_a^a,\quad C_a'=S_aC_a^{a-1},\quad
 C_a^2-S_a^2=1,\quad S_a(0)=0,\quad C_a(0)=1.
\]
For $F\in R$ put $\beta_{a,F}(n)=E(D_a^{2n}F)$, an integer, and
abbreviate $\beta_{a,c^k}$ as $\beta_{a,k}$.
Then $\beta_{a,k}(n)=(2n)![x^{2n}]C_a(x)^k$.
Writing $H_a=\int_0^x C_a^{a-1}$ gives
$C_a\pm S_a=\exp(\pm H_a)$, hence $C_a=\cosh H_a$.
For $m>1$ define the explicit bound
\[
 J(m)=\max\left(\{\lceil p^{2r-1}/2\rceil:
 p^r\parallel m,\ p\text{ odd}\}\cup
 \{2^{r-1}:2^r\parallel m\}\right).
\]
The absent $2$-primary term is omitted when $m$ is odd.

\begin{theorem}\label{thm:allpowers}
For every integer $a\ge1$, $F\in R$ and $m\ge3$,
\[
 \beta_{a,F}\bigl(n+\ph(m)^2/2\bigr)
 \equiv\beta_{a,F}(n)\pmod m\qquad(n\ge J(m)).
\]
Modulo $2$, a period of one holds for $n\ge1$.
The even-exponent specialization for $F=c^k$ is Su's
Theorem 3.2; the argument here also includes every odd exponent
and every integral polynomial observable.
\end{theorem}

\begin{proof}
For an odd prime power $q=p^r$, the ring
$(\Z/q\Z)[s]/(s^q)$ has a unique square root $g$ of $1+s^2$
with $g(0)=1$: solve the coefficients successively, since $2$ is a unit.
The map $c\mapsto g$ sends $D_a$ to $g^a\,d/ds$, because
$g'=s/g$. Both derivatives descend modulo $s^q$, since $q=0$
in the coefficient ring. Choose an integral polynomial lift of $g^a$
with constant coefficient one and of the image of $F$.
Lemma~\ref{lem:operator} gives the derivative-index shift $\ph(q)^2$
after index $p^{2r-1}$. At index $2n$ this is a column shift
$\ph(q)^2/2$ for $n\ge\lceil p^{2r-1}/2\rceil$.

For $q=2^r$ use the free rank-$2q$ ring
\[
 (\Z/q\Z)[s,c]/(s^q,c^2-s^2-1).
\]
Set $A=D_a^2$. Modulo $2$, $A$ is a derivation. If $a$ is even,
\[
 A(s)=0,\qquad A(c)=cw,\qquad w=(1+s^2)^{a-2};
\]
if $a$ is odd,
\[
 A(s)=sw,\qquad A(c)=cw,\qquad w=(1+s^2)^{a-1}.
\]
In each case $A(w)=0$. On a monomial $s^jc^e$ with $e=0,1$,
$A$ acts as $ew$ in the even case and as $(j+e)w$ in the odd case.
Thus $A^{h+1}=w^h A$ for every $h\ge1$.
The exponent of $1+s^2$ in $w$ is even. Consequently, for $r\ge2$,
$w^{2^{r-2}}=1$ modulo $(2,s^{2^r})$ by Frobenius.
For $r=1$ one has $w=1$ modulo $(2,s^2)$.
Take $h=2^{r-2}$ for $r\ge2$ and $h=1$ for $r=1$.
For an integral matrix representing $A$,
$A^{h+1}\equiv A\pmod2$; the commuting-matrix lifting argument gives
\[
 A^{2^{r-1}}(A^{h2^{r-1}}-I)=0\pmod{2^r}.
\]
This gives the column period $2^{2r-3}=\ph(2^r)^2/2$ for $r\ge2$
and period one for $r=1$, after $n\ge2^{r-1}$.
Every local shift divides $\ph(m)^2/2$ for $m\ge3$;
the Chinese remainder theorem proves the result.
\end{proof}

\begin{corollary}\label{cor:companions}
All five conjectured totient-square periods hold, and for $m\ge3$
the half-square shift holds for the following exact source sequences:
\[
\begin{array}{c|ccccc}
 \text{entry}&\mathrm{A281180}&\mathrm{A281181}&\mathrm{A281183}&\mathrm{A281184}&\mathrm{A281440}\\
 \text{value}&\beta_{4,4}(n-1)&\beta_{4,1}(n)&\beta_{4,2}(n)&\beta_{4,3}(n)&\beta_{8,1}(n).
\end{array}
\]
An onset $J(m)+1$ suffices for the first sequence, and $J(m)$
for the other four. For A281180, Theorem~\ref{thm:main} supplies the
additional bound $N(m)$. Modulo $1$ or $2$ a period of one is eventual.
\end{corollary}
\begin{proof}
The identifications in the table follow from the source definitions
\cite{oeis180,oeis181,oeis183,oeis184,oeis440}; for A281180 use
$S_4'=C_4^4$, including its one-index offset. Apply
Theorem~\ref{thm:allpowers}. These five identifications and their
period conclusion are Su's Corollary 3.3.
\end{proof}

\begin{corollary}\label{cor:neighborperiods}
The sequence A281182 has an eventual period dividing $\ph(m)^2$
for $m\ge3$, with onset $2J(m)$, and dividing $2$ modulo $2$.
The sequences A281427, A281428, A281432, A281435, A281436,
A281439, A281443 and A281444 have an eventual half-square shift
for $m\ge3$. Their corresponding exponents $a$ are, respectively,
$5,5,6,7,7,8,9,9$. These period statements for the additional
entries are consequences, not conjectures posed there;
the stronger A281428 assertions are stated below.
\end{corollary}
\begin{proof}
For $v_j=j![x^j](C_4+S_4)$ one has
$v_{2n}=\beta_{4,1}(n)$ and $v_{2n+1}=\beta_{4,4}(n)$.
A shift of $\ph(m)^2$ in $j$ preserves parity and shifts $n$ by
$\ph(m)^2/2$. Apply Theorem~\ref{thm:allpowers}; modulo $2$ use
its period-one column statement. The other entries are $C_a$
even coefficients or $S_a$ odd coefficients, the latter satisfying
$(2n-1)![x^{2n-1}]S_a=\beta_{a,a}(n-1)$.
Their definitions \cite{oeisneighbors} therefore give the claimed periods,
with onset $J(m)$ for the $C_a$ columns and $J(m)+1$ for the $S_a$ columns.
\end{proof}

\begin{corollary}\label{cor:derivative}
Bala's tentative derivative formula in A281183, Formula, revision 60,
signed 22 December 2025, holds: for $g_0=1$ and
$g_{j+1}=(1+z^2)\,d((1+z^2)g_j)/dz$,
$g_{2n}(0)=\mathrm{A281183}(n)$ for every $n\ge0$.
This is Su's Proposition 3.5.
\end{corollary}
\begin{proof}
Put $u(z)=1+z^2$, $\mathcal D=u^2d/dz$ and
$\mathcal Lf=u\,d(uf)/dz$. Then
$\mathcal L=u^{-1}\mathcal D u$ on $\Q[[z]]$ and
$\mathcal L^j1=u^{-1}\mathcal D^ju$.
At $z=0$ this is $j![x^j]u(S_4(x))=j![x^j]C_4(x)^2$.
Taking $j=2n$ proves the stated formula.
\end{proof}

\section{Residue rows and the neighboring formulas}
The five entries also describe geometric rows in particular residue
rectangles and suggest a more general two-dimensional structure.
The following statement gives precise quantifiers; for $a=4$ it is
Su's Proposition 3.4, and the same argument applies to all $a\ge1$.

\begin{corollary}\label{cor:rows}
Fix $m\ge2$ and write $r_* = \max_{p^r\parallel m}r$.
For each odd $p^r\parallel m$ put
\[
 h_{p^r}=\frac{(p-1)p^{r-1}}2,\qquad
 \alpha_{a,p}=[s^{p-1}](1+s^2)^{-a/2}\pmod p.
\]
Set $h_2=h_4=1$ and $h_{2^r}=2^{r-2}$ for $r\ge3$.
Let $H=\operatorname{lcm}_{p^r\parallel m}h_{p^r}$, and choose
$\Gamma\in\Z/m\Z$ reducing to
$\widetilde\alpha_{a,p}^{\,H/h_{p^r}}$ modulo each odd $p^r$,
where $\widetilde\alpha_{a,p}$ is any integral lift of $\alpha_{a,p}$,
and to one modulo the $2$-primary factor.
For every $F\in R$ and $n\ge r_*$,
\[
 \sum_{t=0}^{r_*}(-1)^{r_*-t}\binom{r_*}{t}
 \Gamma^{r_*-t}\beta_{a,F}(n+tH)\equiv0\pmod m.
\]
In particular, for squarefree $m$ the rows satisfy
$\beta_{a,F}(n+H)=\Gamma\beta_{a,F}(n)\pmod m$ for $n\ge1$.
The scalar $\Gamma$ may be zero at some primes.
\end{corollary}
\begin{proof}
On $\F_p[s]/(s^{p^r})$ the odd-prime square-root representation gives
$D_a^p=\lambda(s^p)D_a$: the same derivation and commutation argument
as in Lemma~\ref{lem:operator} says that $\lambda'=0$.
Its constant is $\alpha_{a,p}$. To see the coefficient formula,
work first modulo $s^p$, write $P^{-1}=\sum b_js^j$, and set
$y=\sum_{j=0}^{p-2}b_js^{j+1}/(j+1)$.
Then $D_a y=1-b_{p-1}y^{p-1}$ modulo $y^p$.
Wilson's theorem gives $(D_a^p y)(0)=b_{p-1}$, identifying the constant.
Frobenius and $s^{p^r}=0$ now give
\[
 D_a^{1+(p-1)p^{r-1}}=\alpha_{a,p}D_a\pmod p.
\]
Thus $(D_a^{2H}-\Gamma)D_a$ is divisible by $p$ on the
finite free module; polynomial factorization handles $H/h_{p^r}$.
Commutativity implies
$(D_a^{2H}-\Gamma)^{r_*}D_a^{r_*}=0\pmod{p^r}$.
At the prime two, the proof of Theorem~\ref{thm:allpowers} gives
$(A^H-1)A=0\pmod2$, where $A=D_a^2$;
hence $(A^H-1)^{r_*}A^{r_*}=0\pmod{2^r}$.
For $n\ge r_*$, multiply by the remaining powers, apply $E$,
and use the Chinese remainder theorem. Expanding the power
of the shift difference proves the formula.
\end{proof}

For $a=4$ one has $\alpha_{4,p}=(-1)^{(p-1)/2}/2$.
For the five rectangles in the sources, the column width and row
multiplier are $(6,2)$ for A281180 modulo $7$, $(6,7)$ for A281181
modulo $13$, $(5,5)$ for A281183 and A281184 modulo $11$, and
$(6,6)$ for A281440 modulo $13$. The finite criterion below also
establishes these exact row identities at their displayed initial indices.
For A281440 the rectangle is read from $n=1$; the sequence with
its stated offset $n=0$ begins $1,1,3,5,7,8,6,6$ modulo $13$.

\begin{corollary}\label{cor:formulas}
The question in A281436, Formula, revision 10, signed
23 December 2025, has answer yes: if
$f_0(\theta)=\sec\theta$ and
$f_{j+1}=\sec^5\theta\,f_j'$, then
$f_{2n}(0)=\mathrm{A281436}(n)$.
In A281428, Formula, revision 54, signed 22 January 2026,
the formula for ordinary coefficients requires the normalization
\[
 f_k(j,0)=j![x^j]C_5(x)^k,
\]
where $f_k(0,\theta)=\sec^{k-3}\theta$ and
$f_k(j+1,\theta)=\partial_\theta(\sec^3\theta f_k(j,\theta))$.
Su's Section 4, Source-formula normalization, already gives this correction.
For A281427 and A281435, Formula, signed 22 December 2025,
the displayed nested-derivative rule with $f_0=1$ instead has
$a(n)=f_{2n-2}(0)$, not $f_{2n}(0)$, with the source offset $n\ge1$.
\end{corollary}
\begin{proof}
Let $\theta'=\sec^{a-2}\theta$, $\theta(0)=0$.
Then $S_a=\tan\theta$ and $C_a=\sec\theta$ satisfy the defining flow.
Repeated differentiation of $\sec\theta(x)$ proves the first assertion
at $a=7$.
For $a=5$ put $P=\sec^3\theta$, $T=P\partial_\theta$ and
$L=\partial_\theta(P\,\cdot)=P^{-1}TP$.
Thus $L^j(\sec^{k-3}\theta)=P^{-1}T^j\sec^k\theta$.
Evaluation at zero proves the factorial formula. In particular,
at $k=1,j=2$ its left side is $1$, whereas $[x^2]C_5=1/2$.
Finally, in the $s$ coordinate put $P=(1+s^2)^{a/2}$,
$T=P\partial_s$ and $L=\partial_s(P\,\cdot)=P^{-1}TP$.
Then $(L^j1)(0)=(T^jP)(0)=S_a^{(j+1)}(0)$.
Take $j=2n-2$. At $n=1$, the unshifted source expressions
$f_2(0)$ equal $5$ and $7$, respectively, whereas each $a(1)=1$.
All identities are formal power-series identities over $\Q$.
\end{proof}

\begin{corollary}\label{cor:inverse}
For the same quadratic flow, put $T_a=S_a/C_a$ and
$\theta_a=\arctan S_a$. The common inverse-function formulas in
the source entries follow from
\[
 \begin{aligned}
 x&=\int_0^{S_a}(1+t^2)^{-a/2}\,dt
   =\int_0^{\theta_a}\cos^{a-2}t\,dt,\\
 x&=\int_0^{H_a}\operatorname{sech}^{a-1}t\,dt
   =\int_0^{T_a}(1-t^2)^{(a-3)/2}\,dt.
 \end{aligned}
\]
In particular $T_a'=C_a^{a-3}$ and $T_a^2+C_a^{-2}=1$.
For $a=5$, $T_5^3-3T_5+3x=0$ and
$(4-9x^2)C_5^6-3C_5^2-1=0$.
For $a=7$, $3T_7^5-10T_7^3+15T_7-15x=0$ and
$(225x^2-64)C_7^{10}+40C_7^4+15C_7^2+9=0$.
These are the algebraic identities recorded in A281427, A281428,
A281435 and A281436. Their $S_a$ identities follow by
$C_a^2=1+S_a^2$.
For $a=4$, the coefficient expansion is
$T_4=x+x^3/3!+13x^5/5!+493x^7/7!+\cdots$;
the denominator $6!$ on the $x^5$ term in A281181, Formula,
14 December 2025, must be $5!$.
\end{corollary}
\begin{proof}
Differentiate each integral using the defining flow; its derivative is
one and its constant is zero. For the last integral,
$T_a'=C_a^{a-3}$ follows by the quotient rule and $C_a^2-S_a^2=1$.
At $a=5,7$ the last integrand is a polynomial, giving the displayed
$T_a$ identities. Substitute $T_a=S_a/C_a$ and square:
$3xC_5^3=S_5(2C_5^2+1)$ and
$15xC_7^5=S_7(8C_7^4+4C_7^2+3)$ give the two $C_a$ equations.
Finally $T_4'=C_4$ gives the corrected factorial by integration.
\end{proof}

\begin{corollary}\label{cor:triangle}
For $u=1+s^2$ and $D=u^2\partial_s$, write
\[
 D^{2n-1}s=u^{2n}Q_n(s^2)\qquad(n\ge1).
\]
Then $Q_n(z)=\sum_k T(n,k)z^k$ is the triangle A394012.
Its first column is A281180 and its second column is A391180, with
\[
 T(n,1)=\tfrac12 a(n+1)-2n a(n).
\]
Consequently, for $m\ge2$, A391180 has the eventual period
$\operatorname{lcm}(m,\ph(2m)^2/2)$, with sufficient onset $N(2m)$.
\end{corollary}
\begin{proof}
Here $Q_1=1$. Applying $D^2$ and setting $h=2n-2$ gives
\[
 Q_n=2h(1+(2h+3)z)Q_{n-1}
 +2(1+z)(1+(4h+5)z)Q_{n-1}'
 +4z(1+z)^2Q_{n-1}''.
\]
The coefficient of $z^k$ is precisely the three-term row rule
in A394012 \cite{oeistriangle}; this proves the identity by induction.
Its constant coefficient is $(D^{2n-1}s)(0)=a(n)$.
Since $D^2g(0)=g''(0)$, the coefficient of $s^2$ in
$D^{2n-1}s$ is $a(n+1)/2$. Multiplication by $u^{-2n}$
subtracts $2n a(n)$, giving the second-column formula.
It equals the exponential coefficient of
$S_4''/2-xS_4'-S_4$, the definition in A391180.
Apply Theorem~\ref{thm:main} modulo $2m$ and shift $n$ by a
multiple of both $m$ and $\ph(2m)^2/2$; both terms in the
second-column identity are then unchanged modulo $m$.
\end{proof}

The two distinct modular conjectures in A281428, Comments, revision 54,
signed 3 March 2026, have the stronger answer supplied by
Su's Theorem 4.1: its terms eventually vanish modulo every modulus
coprime to $6$, and have an eventual period dividing $\ph(m)$ for
every positive modulus $m$. These are prior results; the half-square
consequence of Theorem~\ref{thm:allpowers} does not supply the
stronger vanishing or totient-period conclusions.

\section{Exact certificates for the source examples}
Su's Proposition 3.6 gives the periods and the congruence modulo
$12$ below. Here is a small finite-dimensional criterion that
allows the displayed values to be checked independently.
Write $z=s^2$, $e\in\{0,1\}$ and represent a parity component by
$c^e f(z)$. In this section $a=2d$ and
$b_{d,k}=\beta_{2d,k}$. If $G=ef+2(1+z)f'$, direct differentiation gives
\begin{equation}\label{eq:columnoperator}
 D^2(c^ef)=c^e\mathcal B_{d,e}f,\qquad
 \mathcal B_{d,e}f=(1+z)^{2d-2}
 \left((1+(e+2d-1)z)G+2z(1+z)G'\right).
\end{equation}
For a modulus $m\ge2$ take $h=\lceil m/2\rceil$.
The coefficient of $z^{h-1}$ in $\mathcal B_{d,e}(z^h)$ is
$2h(2h-1)$, divisible by $m$; there are no terms of lower degree.
The ideal $(z^h)$ is therefore stable, also for all its higher
monomials, since the operator lowers degree by at most one.
Thus \eqref{eq:columnoperator} defines an $h\times h$ matrix $B$
modulo $m$, and
\[
 b_{d,k}(n)=[z^0]B^n(1+z)^{(k-e)/2},\qquad e\equiv k\pmod2.
\]
By Cayley--Hamilton, the observations and their shift differences
satisfy a monic recurrence of order $h$ over $\Z/m\Z$.
It suffices to check $h$ consecutive zero differences to prove
the proposed period for the entire tail. This is a finite proof
criterion, not extrapolation from a long periodic prefix.

\begin{corollary}\label{cor:certificates}
The source examples have the following exact least periods.
\begin{center}
\begin{tabular}{@{}lrrrr@{}}\toprule
Entry & $m$ & least period & first index & $h$\\\midrule
A281180 & 7 & 18 & 1 & 4\\
A281181 & 10 & 8 & 0 & 5\\
A281181 & 9 & 3 & 1 & 5\\
A281181 & 13 & 72 & 0 & 7\\
A281183 & 11 & 25 & 1 & 6\\
A281184 & 11 & 25 & 1 & 6\\
A281440 & 13 & 72 & 0 & 7\\\bottomrule
\end{tabular}
\end{center}
Also A281181$(n)=1\pmod{12}$ for all $n\ge0$, and
A281184$(n)=3\pmod9$ for all $n\ge1$.
\end{corollary}
\begin{proof}
The program \texttt{verify.py} gives the matrices, the $h$ zero
differences, and a nonzero witness excluding each shift $P/\ell$
for every prime divisor $\ell$ of the proposed period $P$.
These witnesses prove minimality: a smaller period would have
a common-divisor period with $P$ and therefore a period dividing
one of these $P/\ell$. For example, for A281180 modulo $7$,
$a(1)=1$, $a(10)=6$ and $a(7)=2$ exclude shifts $9$ and $6$.
The same rank-six criterion modulo $12$ gives
$b_{2,1}(n)=1\pmod{12}$ for every $n\ge0$: its values for
$0\le n\le6$ are all one, so six successive differences vanish.
This also proves Curtz's divisibility assertion in A281181.
For A281184 modulo $9$, the six values at indices $1$ through $6$
are all $3$; the rank-five criterion propagates their five successive
zero differences. This proves the constant-tail comparison
cross-referenced in A281181.
\end{proof}

\section{Coverage and formal statements}
The following coverage distinguishes the period questions from
the further material cross-referenced in the source entries.
{\small
\begin{longtable}{@{}p{.30\linewidth}p{.65\linewidth}@{}}\toprule
Source item & Coverage\\\midrule
\endfirsthead
\toprule Source item & Coverage\\\midrule
\endhead
\bottomrule\endfoot
A281180, Comments, revision 86, 18 Feb 2026
 & Reproved; prior proof Su (2026), Corollary 3.3. Theorem~\ref{thm:main}.\\
A281181, Comments, revision 103, 16 Feb 2026
 & Reproved; prior proof Su (2026), Corollary 3.3.
 Corollary~\ref{cor:companions}.\\
A281183, Comments, revision 60; A281184, Comments, revision 41;
 both 16 Feb 2026
 & Reproved; prior proof Su (2026), Corollary 3.3.
 Corollary~\ref{cor:companions}.\\
A281440, Comments, revision 17, 19 Feb 2026
 & Reproved; prior proof Su (2026), Corollary 3.3.
 Corollary~\ref{cor:companions}.\\
Five numerical-period examples; A281181 mod $10$, mod $9$, and difference mod $12$
 & Reproved; prior proof Su (2026), Proposition 3.6.
 Corollary~\ref{cor:certificates}, including the three observations of
 Paul Curtz, 20 Jan 2017, in A281181 Comments.\\
Tentative residue-row structure, five Examples sections
 & Reproved and extended; prior precise formulation Su (2026),
 Proposition 3.4 for $a=4$. Corollary~\ref{cor:rows} covers every $a\ge1$.
 The informal source does not specify a unique universal rectangle.\\
A281183, Formula, 22 Dec 2025
 & Proved; prior proof Su (2026), Proposition 3.5.
 Corollary~\ref{cor:derivative}.\\
Common integral, exponential and inverse-function formulas;
the odd-exponent algebraic formulas
 & Proved consequences of the source flow; Corollary~\ref{cor:inverse}.
 No originality is assigned to the recorded identities.\\
A281181, Formula, 14 Dec 2025, $x^5$ term of $T$
 & Refuted-and-corrected factorial: $13x^5/5!$, rather than
 $13x^5/6!$, since $T_4'=C_4$; Corollary~\ref{cor:inverse}.\\
A281182, revision 32, defining e.g.f.
 & Not posed; proved: the interleaving period in
 Corollary~\ref{cor:neighborperiods}.\\
A281436, Formula, revision 10, 23 Dec 2025: ``Then a(n) = f(2*n, 0)?''
 & Answered: yes. Corollary~\ref{cor:formulas} proves the exact
 identity for all $n\ge0$.\\
A281428, Formula, revision 54, 22 Jan 2026: ``f\_k(n, 0) = [x\^{}n] C(x)\^{}k''
 & Refuted-and-corrected normalization; prior correction Su (2026),
 Section 4. The factor $n!$ is required; witness $(k,n)=(1,2)$.
 Corollary~\ref{cor:formulas}.\\
A281427 and A281435, Formula, 22 Dec 2025: ``a(n) = f(2*n, 0)''
 & Refuted-and-corrected offset: use $f(2n-2,0)$ for $n\ge1$.
 Witness $n=1$, where the displayed expressions give $5$ and $7$
 instead of $1$; Corollary~\ref{cor:formulas}.\\
A394012, revision 47, and A391180, revision 71
 & Proved: derivative-polynomial triangle identity and both first
 columns, Corollary~\ref{cor:triangle}. The second-column eventual period
 is not posed in these sources; it is proved here.\\
A281427, A281428, A281432, A281435, A281436, A281439, A281443, A281444
 & Not posed; proved: the half-square congruences in
 Corollary~\ref{cor:neighborperiods}. Even-exponent instances also follow
 from Su (2026), Theorem 3.2.\\
A281428, Comments, revision 54, 3 Mar 2026, conjectures 1 and 2
 & Prior proof cited: Su (2026), Theorem 4.1 proves eventual
 vanishing away from $6$ and the stronger totient-period bound.
 These stronger statements are not deduced from the present half-square theorem.\\
A281184 mod $9$, compared with the constant-tail entry A122553
 & Proved by the rank-five certificate in Corollary~\ref{cor:certificates}.\\
Other columns of A394012; the asymptotic formula in A281428;
unrelated cross-references
 & Not covered. The triangle identification alone makes no modular
 assertion for columns $k\ge2$.\\
\end{longtable}}

The following pinned Lean files give the formal statements corresponding
to the original five totient-square assertions. All links specify
the same source commit.
\begin{itemize}
\item \href{https://github.com/the-omega-institute/trureturing/blob/92de86888730265f84e24d28bdb67546c543a966/D5/S1/Recurrence/Periodic/AutonomousDerivationTotientSquarePeriod.lean}{Polynomial-flow file}, namespace \texttt{BalaTruncated}:
\texttt{prime\_power\_derivative\_period} states
\eqref{eq:operator} for arbitrary $P(0)=1$ and polynomial observables;
\texttt{integer\_matrix\_totient\_square} states the
integer-matrix divisibility underlying \eqref{eq:lift}.
\item \href{https://github.com/the-omega-institute/trureturing/blob/92de86888730265f84e24d28bdb67546c543a966/D5/S1/Recurrence/Periodic/BalaA281180TotientSquarePeriod.lean}{A281180 source}:
\texttt{result} proves integrality and the eventual shift $\ph(m)^2$
for the source equation \eqref{eq:source}.
\item The declarations \texttt{result} in the
\href{https://github.com/the-omega-institute/trureturing/blob/92de86888730265f84e24d28bdb67546c543a966/D5/S1/Recurrence/Periodic/BalaA281181TotientSquarePeriod.lean}{A281181},
\href{https://github.com/the-omega-institute/trureturing/blob/92de86888730265f84e24d28bdb67546c543a966/D5/S1/Recurrence/Periodic/BalaA281183TotientSquarePeriod.lean}{A281183},
\href{https://github.com/the-omega-institute/trureturing/blob/92de86888730265f84e24d28bdb67546c543a966/D5/S1/Recurrence/Periodic/BalaA281184TotientSquarePeriod.lean}{A281184} and
\href{https://github.com/the-omega-institute/trureturing/blob/92de86888730265f84e24d28bdb67546c543a966/D5/S1/Recurrence/Periodic/BalaA281440TotientSquarePeriod.lean}{A281440 source files}
prove integrality and the eventual shift $\ph(m)^2$ for $C_4$,
$C_4^2$, $C_4^3$ and $C_8$, respectively, with their hyperbolic
defining equations. The A281184 declaration also includes
$(\log(C_4+S_4))'=C_4^3$.
\end{itemize}
The all-exponent extension, the half-square shifts, explicit starting
indices, residue-row extension, interleaving, neighboring derivative and
triangle identities, and finite-certificate interpretation
are supplied by the written arguments above. The exact numerical
certificates and coefficient comparisons are provided by
\texttt{verify.py} and \texttt{checks.json}; they are separate from
these Lean declarations. The coefficient comparisons cover 21
terms in each family, the first eight against the OEIS values,
and derivative indices $0$ through $40$ in \eqref{eq:source}.
The polynomial-operator checks cover five polynomials at each of twelve
prime-power moduli, hence 60 finite cases. A further 108 operator
identities use the full quadratic algebra at exponents $1$ through $9$
and the same twelve prime powers. Five residue-row certificates and
the additional constant-tail certificate use the finite criterion above.
The triangle and derivative comparisons have their exact finite ranges
in \texttt{checks.json}; their all-index scope follows from the written proofs.

\begin{thebibliography}{9}\small
\bibitem{su} L.-H. Su, \emph{Quadratic Derivations and Uniform Modular
Periods for Six Inverse-Function Sequences}, version 1, 31 October
2026, Theorem 3.2, Corollary 3.3 and Propositions 3.4--3.6.
\href{https://doi.org/10.5281/zenodo.23835577}{doi:10.5281/zenodo.23835577}.
\bibitem{oeis180} OEIS Foundation Inc., \emph{A281180}, revision 86,
6 November 2026; conjecture by P. Bala, 18 February 2026.
\url{https://oeis.org/A281180}.
\bibitem{oeis181} OEIS Foundation Inc., \emph{A281181}, revision 103,
6 November 2026; conjecture by P. Bala, 16 February 2026.
\url{https://oeis.org/A281181}.
\bibitem{oeis183} OEIS Foundation Inc., \emph{A281183}, revision 60,
6 November 2026; conjecture by P. Bala, 16 February 2026.
\url{https://oeis.org/A281183}.
\bibitem{oeis184} OEIS Foundation Inc., \emph{A281184}, revision 41,
6 November 2026; conjecture by P. Bala, 16 February 2026.
\url{https://oeis.org/A281184}.
\bibitem{oeis440} OEIS Foundation Inc., \emph{A281440}, revision 17,
6 November 2026; conjecture by P. Bala, 19 February 2026.
\url{https://oeis.org/A281440}.
\bibitem{bala} P. Bala, \emph{The method of Graves for computing
inverse functions}, 26 January 2026, Theorem 1 and Corollary 1.
\url{https://oeis.org/A281181/a281181.pdf}.
\bibitem{dominici} D. Dominici, \emph{Nested derivatives: a simple method
for computing series expansions of inverse functions}, Int. J. Math.
Math. Sci. \textbf{2003}, no. 58 (2003), 3699--3715, Theorem 4.1.
\href{https://doi.org/10.1155/S0161171203303291}{doi:10.1155/S0161171203303291};
\href{https://arxiv.org/abs/math/0501052}{arXiv:math/0501052}, version 2 (2005).
\bibitem{danesh} P. Danesh, \emph{Eventual periodicity of weighted
Stirling transforms and two conjectures of Bala}, Cambridge Open Engage,
version 1, 26 August 2026, Theorems 3.1 and 4.1.
\href{https://doi.org/10.33774/coe-2026-4d6jr}{doi:10.33774/coe-2026-4d6jr}.
\bibitem{oeisneighbors} OEIS Foundation Inc., neighboring inverse-function entries,
versions accessed 14 November 2026:
\href{https://oeis.org/A281182}{A281182},
\href{https://oeis.org/A281427}{A281427},
\href{https://oeis.org/A281428}{A281428},
\href{https://oeis.org/A281432}{A281432},
\href{https://oeis.org/A281435}{A281435},
\href{https://oeis.org/A281436}{A281436},
\href{https://oeis.org/A281439}{A281439},
\href{https://oeis.org/A281443}{A281443},
\href{https://oeis.org/A281444}{A281444}.
The A281436 question is signed P. Bala, 23 December 2025;
the A281428 modular conjectures are signed P. Bala, 3 March 2026.
\bibitem{oeistriangle} S. Manyama, OEIS
\href{https://oeis.org/A394012}{A394012}, revision 47,
and \href{https://oeis.org/A391180}{A391180}, revision 71,
13--14 April 2026; defining recurrence and second-column formula.
\end{thebibliography}
\end{document}
